\section{Mapping Algebra}%
\label{sec:mapping_algebra}

% {
% \color{blue}
%     In this section, we define the algebraic mapping operators used
%     in our formalism of mapping processes to generate RDF datasets 
%     form (semi-)structured data. 
%     For each operator, we also introduce the rationale behind the 
%     choice of defining such operators for use in our formalism. 
%     Throughout this section, a running example is provided by 
%     applying the operators sequentially on a sample dataset. 
%     First, we define the \emph{mapping definition} as follows: 
% \begin{definition}
%     \label{def:mapping_definition}
%     A \textbf{mapping definition}  is an expression that can be built by 
%     combining various operators of the mapping algebra. 
%     It represents the logical plan of generating RDF datasets from  
%     (semi-)structured data.
% \end{definition}
% }

This section introduces our mapping algebra which consists of six
	\todo{Change this number if we decide to remove some of the operators due to space constraints.}
types of operators.
	The output of each such operator is
	%The output of each operator is
	%Each such operator outputs
a mapping relation and the inputs (if any) are mapping relations as well. Hence, the operators can be combined
	%to form
	into
arbitrarily complex algebra expressions, where each such expression defines a mapping from one or more sources of
	%structured or semi-structured
	(semi-)structured
data into an~RDF~dataset.\footnote{We assume here that the mapping relation resulting from the root operator of the expression contains the aforementioned attributes $\subjAttr$, $\predAttr$, $\objAttr$, and $\graphAttr$ (see Definition~\ref{def:generated_dataset}).}


\subsection{Source Operator}%
\label{sub:source_operator}

% { \color{blue}
% The mapping relations (Definition~\ref{def:mapping_relation}), which serve as
% the data model for the algebraic operators, must first be generated from
% (semi-)structured data. 
% To accomplish this, a dedicated operator, the Source
% operator, is introduced to produce mapping relations from (semi-)structured
% data. 
% Before defining the Source operator, we present essential notations and
% concepts involved in its definition.
% }
% 
% \noindent \textbf{Notations}  We assume the following pairwise disjoint countably infinite sets:
% $\symDatasetUni$ (datasets), and  $\symDataObjUni$ (data objects).
% Each dataset in $\symDatasetUni$ is a collection of data objects.
% To denote that a data object~$\dataObject \in \symDataObjUni$ is contained in
% such a dataset~$\symDataset \in \symDatasetUni$, we write
% $\dataObject \in \symDataset$.

The basis of every
	expression that defines a mapping using our algebra
	%algebraic mapping expression
are operators that establish a relational view of each source of input data. From the perspective of the algebra, these \emph{source operators} are nullary operators; they do not have any mapping relation as input. The mapping relation that such an operator creates
	as output
depends on the data of the corresponding data source and can be specified based on a language appropriate to identify and extract relevant pieces of data from the data source. As our mapping algebra should be applicable for various types of data sources, for which different such languages are suitable, we introduce an abstraction of this extraction process and define
	our notion of source operators
	%the source operator
based on this abstraction.

	%For our abstraction we assume
	Our abstraction considers
two infinite sets, $\symDataObjUni$ and $\symDataAccUni$, which are disjoint.
	\todo{Min Oo, notice that I have removed the set $\symDatasetUni$ of datasets, which we said are collections of data objects. I don't think we need that. After all, a dataset can also be seen as a data object; especially because data objects may contain other data objects. As a consequence, the purpose of having an explicit notion of datasets, distinct from data objects, is not easy to explain. Therefore, I have removed it. Once this change is implemented throughout the rest of the paper, I would prefer to change the symbol $\symDataObjUni$ to $\mathcal{D}$.}
%
$\symDataObjUni$ is the set of all possible \emph{data objects}.
	%In concrete instantiations of our abstraction,
	In practice,
there are different kinds of data objects, even within the context of the same type of data source, and data objects may be nested within one another. For instance,
	%in the context of JSON documents, each JSON object may be seen as a data object, and so may each named field within a JSON object and also the name and the value of such a field.
	each line of a CSV file may be seen as a data object, and so may each value within such a line, as well as the CSV file as a whole.
In our abstraction, any particular kind of data objects is considered as a specific subset of $\symDataObjUni$, and to indicate that a data object~$\dataObject \in \symDataObjUni$ is contained in another data object~$\symDataset \in \symDataObjUni$, we write $\dataObject \in \symDataset$.

The set $\symDataAccUni$ captures possible \emph{query languages} to retrieve data objects from other data objects. We consider each such language $\symDataAcc \in \symDataAccUni$ as a set, where every element $\queryExpr \in \symDataAcc$ is one of the expressions
	admitted by
	%of
the language. In practice, each such language is designed to be used for a specific kind of data objects (i.e., a specific subset of $\symDataObjUni$).
	%Some languages may also be repurposed to be used for different kinds of data objects.
Therefore, we define the notion of the \emph{types of data sources} considered for the source operators of our algebra as a structure that specifies the relevant kind of data objects and the corresponding evaluation semantics of the language to be used for data sources of such a type.

% Informally, a data object~$\dataObject \in \symDataObjUni$ is an abstract
% entity that could be queried using a query expression to extract
% one or more data objects \textbf{which can, in turn, be accessible by different data access languages}.
% For example, a JSON object is considered a data object with JSONPath as the data
% access language, and it can be further queried using JSONPath expressions to
% extract JSON arrays (data objects) or a JSON field containing a string
% representation of CSV rows (data objects), which can be further queried using a CSV specific
% data access language.

\smallskip\noindent
\textbf{\color{red} Everything below is still work in progress (I am still thinking about putting $\cbeval$ back)}

\begin{definition}
	\label{def:source_type}
	\normalfont
	A \textbf{source type} is a tuple $\sourceTypeTuple$ in which
	$\symSetDataset \subseteq \symDataObjUni$,
	$\symDataAcc$ is a query language in $\symDataAccUni$,
	and $\eval$ is a partial function $\eval \!: \symSetDataset \times \symDataObjUni \times \symDataAcc \rightarrow \symDataSeqUni$,
\todo{Min Oo, notice that I have decided to replace the two evaluation functions by a single one. When I tried to describe their purpose, I realized that there is not really a totally convincing argument for having two separate ones. \\ Attention: this change still needs to be reflected in the rest of the paper.}
	where $\symDataSeqUni$ is the set of all possible sequences of data objects in $\symDataObjUni$.
	%
	A \textbf{data source}~$\dataSource$ is a tuple~$\dataSourceTuple$ in which
	$\sourceType$ is a source type $\sourceTypeTuple$ and
	$\symDataset \in \symSetDataset\!$.
\end{definition}

As per Definition~\ref{def:source_type}, the set $\symSetDataset$ of a source type~$\sourceTypeTuple$ determines the kind of data objects that the data sources of this type may provide access to.
	For instance, for CSV-based data sources this may be the (infinite) set of all possible CSV files. %, and for JSON-based data sources, it may be the (infinite) set of all JSON documents.
	%For instance, for $\sourceTypeX{\mathsf{csv}} = \sourceTypeTupleX{\mathsf{csv}}$, which represents the type of all CSV-based data sources, $\symSetDatasetX{\mathsf{csv}}$ is the (infinite) set of all possible CSV files.
The function~$\eval$
	is meant to define
	%captures
the evaluation semantics of the language~$\symDataAcc$
	when used
for data sources of the given type%
	%~(which opens the option to repurpose the same language for different types of data sources)%
. The domain of this function consists of 3-tuples that contain an additional data object as their second element. This is to cover cases in which the evaluation of a query $\queryExpr \in \symDataAcc$ assumes a context object, in addition to the dataset $\symDataset \in \symSetDataset$ over which
	%the query
	$\queryExpr$
is evaluated. JSONPath and XPath are typical examples of languages that assume a context object. For languages for which this is not the case, the second argument of the $\eval$ function may simply be fixed to be the same as the first~argument.

	While a formal definition of concrete source types is beyond the scope of this paper, the
	%The
following two examples outline informally how JSON-based data sources and CSV-based data sources may be captured within our abstraction.

\begin{example}
	\label{ex:source_type_json}
	We may introduce $\sourceTypeX{\mathsf{json}} = \sourceTypeTupleXXX{\mathsf{json}}{\mathsf{jp}}{\mathsf{jp}}$ as a source type for JSON-based data sources. To this end, $\symSetDatasetX{\mathsf{json}}$ is the set of all possible JSON documents, $\symDataAccX{\mathsf{jp}}$ is the set of all
		possible
	JSONPath expressions, and $\evalX{\mathsf{jp}}$ captures the evaluation semantics of JSONPath
		%expressions over JSON documents.
		expressions.
	For instance, for a JSON document $\symDataset \in \symSetDatasetX{\mathsf{json}}$ and a JSONPath expression $\queryExpr \in \symDataAccX{\mathsf{jp}}$ that begins with the selector for the root element (i.e., \texttt{\$}), $\evalX{\mathsf{jp}}(\symDataset,\symDataset,\queryExpr)$ returns the elements of $\symDataset$ selected by $\queryExpr$. If $\queryExpr'\! \in \symDataAccX{\mathsf{jp}}$ is a JSONPath expression that does not begin with~\texttt{\$}, and assume $\dataObject$ is a JSON object within $\symDataset$, then $\evalX{\mathsf{jp}}(\symDataset,\dataObject,\queryExpr')$ returns the elements of $\symDataset$ that can be reached from $\dataObject$ as per the traversal specified by $\queryExpr'\!$.
\end{example}

\begin{example}
	\label{ex:source_type_csv}
	To define $\sourceTypeX{\mathsf{csv}} = \sourceTypeTupleX{\mathsf{csv}}$ as a source type for CSV-based data sources, we let $\symSetDatasetX{\mathsf{csv}}$ be the infinite set of all possible CSV files. As the language $\symDataAccX{\mathsf{csv}}$, we may use the set of all strings (i.e., $\symDataAccX{\mathsf{csv}} = \symAllStrings$). Informally, $\evalX{\mathsf{csv}}$ may then be defined as follows. For the empty string~$\varepsilon$, and every CSV document~$\symDataset \in \symSetDatasetX{\mathsf{csv}}$, $\evalX{\mathsf{csv}}(\symDataset,\symDataset,\varepsilon)$ returns all rows of $\symDataset$%
		%, in the order in which they appear within $\symDataset$%
	. Moreover, for every CSV document~$\symDataset \in \symSetDatasetX{\mathsf{csv}}$, every string $c$ that is a column name in $\symDataset$, and every row $\dataObject \in \symDataset$, $\evalX{\mathsf{csv}}(\symDataset,\dataObject,c)$ is the value that
		the row~%
	$\dataObject$ has in column~$c$.
\end{example}

%The question of how data sources of any such type are actually accessed in practice is irrelevant for our formalization.

As a last ingredient for
	%the following definition of the
	defining
source operators, we assume a partial function $\typeCast \!: \symDataObjUni \rightarrow \symAllLiterals$ that maps (some) data objects to RDF literals. For instance, string values in CSV files and in JSON documents may be mapped to \ttl{xsd:string} literals, numeric JSON values that are integers may be mapped to \ttl{xsd:integer} literals, etc.

\todo[inline]{Olaf, continue here with a brief informal description of how the source operators work ... }

% \begin{definition}
% 	\label{def:cast_function}
% 	A \textbf{type-cast} function is a partial function
% 	$\typeCast: \symDataObjUni \rightarrow \symAllLiterals$.
% \end{definition}

\begin{definition}
	\label{def:source_operator}
	\normalfont
	Let $\dataSource = \dataSourceTuple$ be a data source with
	$\sourceType=\sourceTypeTuple$,
	$\queryExpr \in \symDataAcc$, and
	$\symAttrQueryMap$ be a partial function
	$\symAttrQueryMap \!: \symAttrUniverse  \rightarrow \symDataAcc$.
	The \textbf{$(\queryExpr,\symAttrQueryMap)$-specific mapping relation obtained from $\dataSource$}, denoted by $\sourceOpDflt$, is the mapping relation $\mappingRel$ such that
	\begin{align*}
		\symAttrSubSet  & = \fctDom{\symAttrQueryMap},
		\; \text{ and}
		\\
		\symMappingInst & = \left\{ \{\attr_1 \rightarrow \fctTypeCast{\var_1}, \dots, \attr_n \rightarrow \fctTypeCast{\var_n} \} \ | \
		\dataObject \in \symDataSeq_{root} \text{ and } ((\attr_1 , \var_1), \dots, (\attr_n , \var_n)) \in X_d\right\}
	\end{align*}
	where $\symDataSeq_{root} = \fctEval{\symDataset}{\queryExpr_{root}}$ and
	\begin{equation}
		X_d = \bigtimes_{\attr \in \fctDom{\symAttrQueryMap}}
		\{(\attr, \var) \ | \
		\var \in \fctCbeval{\symDataset}{\dataObject}{\queryExpr} \text{ with }  q = \symAttrQueryMap(\attr)\}
		.
	\end{equation}
\end{definition}



\begin{example}
\label{ex:source_operator}
We provide a running example used throughout the rest of the paper starting
with a CSV file with the header, $H= [(id, firstname, lastname, age)]$, 
containing the following 2 rows of CSV data records on people as: 
\begin{equation*}
    \symDataset = [(1, \text{Alice}, \text{Lee}, 23), (2, \text{Bob}, \text{Malice}, 25)] 
\end{equation*}
    A source operator, $\sourceOpDflt$, with $\dataSource = \dataSourceTuple$,
    $\sourceType = \sourceTypeTuple$, $\queryExpr_{root}$ a CSV root query 
    which, when evaluated as $\eval(\queryExpr_{root})$, will return the 2 CSV records in $\symDataset$
    as data objects, and $\symAttrQueryMap = \{?id \rightarrow id,  ?name \rightarrow firstname, ?age \rightarrow age\}$    
    generates a mapping relation $r$  with two mapping tuples $t_1, t_2 \in r$ as:
    \begin{align*}
        t_1 &= \{ ?id \rightarrow 1, ?name \rightarrow \text{Alice}, ?age \rightarrow 23 \}\quad t_2 = \{ ?id \rightarrow 2, ?name \rightarrow \text{Bob}, ?age \rightarrow 25 \} \\
    \end{align*}
    Do note that due to the $cast$ function (Definition~\ref{def:cast_function}), 
    all the mapping tuples $t \in r$ only have RDF literals in the $range(t)$.
\end{example}


\subsection{Projection Operator}%
\label{sub:projection_operator}

The \textbf{projection} operator restricts the mapping tuples in a mapping
relation to a subset of the attributes of that relation.


\begin{definition}
	\label{def:tuple_restrict}
	Given a mapping tuple $\mappingTuple$ and a subset $A \subseteq
		\mathrm{dom}(\mappingTuple)$ of the attributes for which $\mappingTuple$ is
	defined, we write $\mappingTuple[A]$ to denote the mapping tuple
	$\mappingTuple'$ that is the restriction of $\mappingTuple$ to $A$;
	i.e., $\mathrm{dom}(\mappingTuple') = A$ and $\mappingTuple'(a) =
		\mappingTuple(a)$ for all $a \in A$.
\end{definition}

\begin{definition}
	\label{def:projection_operator}
    \normalfont
	Let $r = \mappingRel$ be a mapping relation and $\symProjSet \subseteq
		\symAttrSubSet$ be a non-empty subset of the
	attributes of $r$.
	The \textbf{$\symProjSet$-specific projection of $r$}, denoted by
    $\fctProjectOpDflt{r}$, is the mapping relation $(\symAttrSubSet',
		\symMappingInst')$ such that $\symAttrSubSet' = \symProjSet$ 
        and $\symMappingInst' = \{ \mappingTuple[\symAttrSubSet'] \ | \
		\mappingTuple \in \symMappingInst \}$.
\end{definition}

\begin{example}
    \label{ex:projection_operator}
    Let $r$ be the mapping relation generated in Example~\ref{ex:source_operator},
    and $\symProjSet = \{?id, ?name\}$. 
    The mapping relation $r'$ generated after applying $\fctProjectOpDflt{r}$
    has mapping tuples $t'_1, t'_2 \in r'$ as:
    \begin{equation*}
        t'_1 = \{ ?id \rightarrow 1, ?name \rightarrow \text{Alice}\},
        t'_2 = \{ ?id \rightarrow 2, ?name \rightarrow \text{Bob}\}
    \end{equation*}
\end{example}


%\subsection{Rename Operator}%
%\label{sub:rename_operator}
%
%
%\begin{definition}
%	\label{def:rename_operator}
%	Let $r = \mappingRel$ be a mapping relation.
%	Given a partial function
%	$\symRenamePFunc: \symAttrSubSet \rightarrow \symAttrUniverse$
%	that is injective,
%    the \textbf{$\symRenamePFunc$-specific renaming of $r$}, denoted by $\fctRenameOp{r}$, 
%	is the mapping relation $(\symAttrSubSet'\!, \symMappingInst')$ such that
%	\begin{align*}
%		\symAttrSubSet'  & =  (\symAttrSubSet \setminus dom(\symRenamePFunc)) \; \cup \;
%		\{ \symRenamePFunc(a) \ |  \ a \in dom(\symRenamePFunc) \}, \; \text{ and }                                       \\
%		\symMappingInst' & = \{ \fctRenameTuple{\symRenamePFunc}{\mappingTuple}\ | \ \mappingTuple \in \symMappingInst\}
%	\end{align*}
%	where, for every mapping tuple~$\mappingTuple \in \symMappingInst$, we write $\fctRenameTuple{\symRenamePFunc}{\mappingTuple}$ to denote the mapping tuple~$\mappingTuple'$ for which it holds that $\fctDom{\mappingTuple'} = \symAttrSubSet'$ and, for every attribute~$\attr \in \fctDom{\mappingTuple'}$,
%	\begin{equation*}
%		\mappingTuple'(\attr) =
%		\begin{cases}
%			\mappingTuple(\symRenamePFunc^{-1}(a)) & \text{if $\attr \in \mathrm{img}(\symRenamePFunc)$,} \\
%			\mappingTuple(a)                       & \text{else.}
%		\end{cases}
%	\end{equation*}
%\end{definition}

\subsection{Extend Operator}%
\label{sub:extend_operator}

\begin{definition}
	\label{def:ExtensionFunction}
	An \textbf{extension function} is a function $\extFunc: \extFuncType \times \dots \times \extFuncType \rightarrow \extFuncType$.
	We denote the set of all possible \textbf{extension function} as $\symExtFuncUni$.
\end{definition}

\begin{definition}
	An \textbf{extend expression}~$\extExpr$ is an expression that is defined recursively as
	follows:
	\begin{enumerate}
		\item
		    Every RDF term in $\symTermUniverse$ is an \textbf{extend expression}.
		\item
            Every attribute in $\symAttrUniverse$ is an \textbf{extend expression}.
		\item
		      If $\extExpr_1,\dots,\extExpr_n$ are extend expressions and
		      $\extFunc \in \symExtFuncUni$, then the tuple $\extFuncTuple$ is an
		      extend expression.
	\end{enumerate}
	% 	Thus, a grammar that accepts such extend expression $\extExpr$ is defined as
	% 	follows where $x \in \symTermUniverse$, $a \in \symAttrUniverse$, and
	% 	$\extFunc \in \symExtFuncUni$:
	% 	\begin{equation*}
	% 		\extExpr ::= x \ |\ a \ |\  (\extFunc, \extExpr, \dots, \extExpr)
	% 	\end{equation*}
\end{definition}

For every extend expression~$\extExpr$, we write $\fctAttrs{\extExpr}$ to denote the set of all attributes mentioned in $\extExpr$. Formally, this set is defined recursively as follows.
\begin{enumerate}
	\item If $\extExpr$ is an RDF term, then $\fctAttrs{\extExpr} = \emptyset$.
	\item If $\extExpr$ is an attribute $\attr \in \symAttrUniverse$, then $\fctAttrs{\extExpr} = \{\attr\}$.
	\item If $\extExpr$ is of the form $\extFuncTuple$, then $\fctAttrs{\extExpr} = \bigcup_{i \in \{1,\ldots,n\}} \fctAttrs{\extExpr_i}$.
\end{enumerate}


\begin{definition}
	\label{def:extend_operator}
	Let $r = \mappingRel$ be a mapping relation, $\attr \in \symAttrUniverse$ be
	an attribute that is not in $\symAttrSubSet$ (i.e., $\attr \notin \symAttrSubSet$),
	and $\extExpr$ be an extend expression that is valid for $r$.
	The \textbf{$\extExpr$-specific extension of $r$ for $\attr$}, 
    denoted by $\fctExtOp{r}$,
	is the mapping tuple
	$(\symAttrSubSet'\!,\symMappingInst')$ such that
	\begin{align*}
		\symAttrSubSet'  & = \symAttrSubSet \cup \{ \attr \}, \; \text{ and }                                            \\
		\symMappingInst' & = \bigl\{\mappingTuple \cup \{\attr \rightarrow \fctEval{\extExpr}{\mappingTuple}\} \ | \
		\mappingTuple \in \symMappingInst \bigr\}
	\end{align*}
	where, for every mapping tuple~$\mappingTuple \in \symMappingInst$, $\fctEval{\extExpr}{\mappingTuple}$ is either an RDF term or the
	error value ($\error$), which is determined recursively as follows:

	\begin{enumerate}
		\item If $\extExpr$ is an RDF term $x$, then $\fctEval{\extExpr}{\mappingTuple} = x$.
		\item If $\extExpr$ is an attribute $\attr$ and $\attr \in \fctDom{\mappingTuple}$, then $\fctEval{\extExpr}{\mappingTuple} = \mappingTuple(\attr)$.
		\item If $\extExpr$ is an attribute $\attr$ and $\attr \notin \fctDom{\mappingTuple}$, then $\fctEval{\extExpr}{\mappingTuple} = \error$.
		\item If $\extExpr$ is a tuple $\extFuncTuple$ such that the extend function~$\extFunc$ is an $n$-ary function, then
		      $\fctEval{\extExpr}{\mappingTuple} =  \fctExtFuncEvalN{\extExpr}{\mappingTuple}$.
		\item If $\extExpr$ is a tuple $\extFuncTuple$ such that the extend function~$\extFunc$ is not an $n$-ary function, then
		      $\fctEval{\extExpr}{\mappingTuple} =  \error$.
	\end{enumerate}

	\begin{equation*}
		\fctEval{\extExpr}{\mappingTuple} =
		\begin{cases}
			\extExpr                                  & \text{if $\extExpr$ is an RDF term;}                                                        \\
			\mappingTuple(\extExpr)                   & \text{if $\extExpr \in \symAttrUniverse$ and $\extExpr \in \fctDom{\mappingTuple}$;}        \\
			\fctExtFuncEvalN{\extExpr}{\mappingTuple} & \ \parbox{15em}{if $\extExpr = \extFuncTuple$ such that $\extFunc$ is an $n$-ary function;} \\
			\error                                    & \text{else.}
		\end{cases}
	\end{equation*}


\end{definition}

\begin{example}
    \label{ex:extend_operator}
    The mapping relation, $r$, generated in
    Example~\ref{ex:projection_operator} only contains RDF literals, thus, an
    RDF dataset containing RDF graphs with triples cannot be generated from
    $r$. 
    Extend operators, applied in succession, can generate the IRI terms
    required for the triples in the resulting RDF dataset. 
    Let $\extExpr_{iri}^{\attr}=(f, \attr)$ be an extend expression with a function~$f$
    that generates an IRI, when evaluating $\eval(\extExpr_{iri}^{\attr},
    \mappingTuple)$, by prepending the URL ``\texttt{example.com/}'' to
    $\mappingTuple(\attr)$.
    We apply 4 extend operators in succession on $r$ as follows:
    \begin{equation*}
        \fctExtOpXBig{\subjAttr}{\extExpr_{iri}^{?id}}{
            \fctExtOpXBig{\predAttr}{\texttt{foaf:name}}{
                \fctExtOpXBig{\objAttr}{?name}{
                    \fctExtOpXBig{\graphAttr}{\texttt{rr:defaultGraph}}{
                        r
                    }
                }
            }
        }
    \end{equation*}
    This results in a new mapping relation with $r'$ with the following 
    mapping tuples $t_1$, and $t_2$ with both having a mapping
    $\graphAttr \rightarrow \texttt{rr:defaultGraph}$:
    \begin{align*}
        t_1 &= \{\dots,
        \subjAttr \rightarrow \text{<example.com/1>},
        \predAttr \rightarrow \texttt{foaf:name}, 
        \objAttr \rightarrow \text{``Alice''}, 
        \dots
        \}\\
        t_2 &= \{\dots,
        \subjAttr \rightarrow \text{<example.com/2>},
        \predAttr \rightarrow \texttt{foaf:name}, 
        \objAttr \rightarrow \text{``Bob''}, 
        \dots
        \}
    \end{align*}
\end{example} 


\subsection{Binary Operators}%
\label{sub:equi_join_operator}

\begin{definition}
	\label{def:equijoin_operator}
	Let $r_{1} = \mappingRelSpec{1}$ and
	$r_{2} = \mappingRelSpec{2}$ be mapping relations
	such that $\symAttrSubSet_{1} \cap \symAttrSubSet_{2} = \emptyset$,
	and let $\symJoinAttrPairs \subseteq \symAttrSubSet_{1} \times \symAttrSubSet_{2}$.
	The \textbf{$\symJoinAttrPairs$-specific equijoin of $r_{1}$ and $r_{2}$},
	denoted by $\fctEquiJoinOp{r_{1}}{r_{2}}$, is the mapping relation
	$(\symAttrSubSet'\!, \symMappingInst')$
	such that $\symAttrSubSet'\! = \symAttrSubSet_{1} \cup \symAttrSubSet_{2}$
	and
	\begin{equation*}
		\symMappingInst'\! =  \{ \mappingTuple_1 \cup \mappingTuple_2 \ | \
		\mappingTuple_1 \in \symMappingInst_{1} \text { and }
		\mappingTuple_2 \in \symMappingInst_{2} \text{ such that }
		\mappingTuple_1(\attr_1) = \mappingTuple_2(\attr_2) \text{ for all }
		(\attr_1,\attr_2) \in \symJoinAttrPairs
		\}
		.
	\end{equation*}
\end{definition}


\begin{example}
    \label{ex:equi_join_operator}
    Example~\ref{ex:projection_operator} generates a mapping relation $r_1$
    containing only the first name of the people by projecting away the other
    attributes. 
    Assuming there is a mapping relation~$r_2$ derived from 
    $D$ (Example~\ref{ex:source_operator}), which contains the last names
    of the people with their corresponding $id$ number, we can get both 
    the first and last name of the people by applying an equijoin on 
    the $id$ attribute with $\symJoinAttrPairs = \{ (?id, ?id)\}$. 
    The resulting mapping relation generated by applying
    $\fctEquiJoinOpBig{r_1}{r_2}$ contains the following two tuples $t_1$ and
    $t_2$: 
    \begin{align*}
        t_1 &=  \{ ?id \rightarrow 1, ?firstname \rightarrow \text{Alice}, ?lastname \rightarrow \text{Lee} \}\\
        t_2 &=  \{ ?id \rightarrow 2, ?firstname \rightarrow \text{Bob}, ?lastname \rightarrow \text{Malice} \}\\
    \end{align*}
\end{example}



\begin{definition}
	\label{def:natural_join_operator}
	Let $r_{1} = \mappingRelSpec{1}$ and
	$r_{2} = \mappingRelSpec{2}$ be mapping relations.
	The \textbf{natural join of $r_{1}$ and $r_{2}$}, denoted by
	$\fctNatJoin{r_{1}}{r_{2}}$, is the mapping relation $(\symAttrSubSet', \symMappingInst')$
	such that $\symAttrSubSet'\! = \symAttrSubSet_{1} \cup \symAttrSubSet_{2}$
	and
	\begin{equation*}
		\symMappingInst' =  \{ \mappingTuple_1 \cup \mappingTuple_2 \ | \
		\mappingTuple_1 \in \symMappingInst_{1} \text { and }
		\mappingTuple_2 \in \symMappingInst_{2} \text{ such that }
		\mappingTuple_1(\attr) = \mappingTuple_2(\attr) \text{ for all }
		\attr \in \symAttrSubSet_1 \cap \symAttrSubSet_2
		\}.
	\end{equation*}
\end{definition}


\begin{definition}
	\label{def:union_operator}
	Let $r_{1} = \mappingRelSpec{1}$ and
	$r_{2} = \mappingRelSpec{2}$ be mapping relations
    such that $\symAttrSubSet_1 = \symAttrSubSet_2$.
	The \textbf{union of $r_{1}$ and $r_{2}$}, denoted by
	$\fctUnion{r_{1}}{r_{2}}$, is the mapping relation $(\symAttrSubSet_1, \symMappingInst_1 \cup \symMappingInst_2)$.
\end{definition}




